\section{Conclusion}
\label{sec:conclusion}

This study examined commodity tail dependence and exchange-rate risk in Ghana using two co-equal calendar constructions, explicit marginal-model adequacy checks, eight static bivariate copula families, finite-threshold tail-concentration measures and multiplicity-controlled stress inference.

Three conclusions remain defensible after those checks. First, no stress--calm comparison survives either Bonferroni or Benjamini--Hochberg correction under the Combined, COVID-19 or 2024 stress definitions in either calendar panel. Second, Gold--Brent, Gold--WTI and Brent--WTI reject all eight tested static copula families in both panels, while Cocoa--Gold, Cocoa--Brent and Cocoa--WTI are calendar-sensitive. Third, the cedi remains inadequately represented by the tested continuous marginal candidates in both panels, so commodity--cedi copula and EVT quantities should be interpreted as conditional diagnostics.

The broader implication is that robustness in dependence analysis cannot be inferred from one preprocessing rule, one selected copula or one nominally significant cell. For cross-market risk analysis, calendar construction, marginal adequacy, absolute goodness of fit and multiple-testing control should be treated as part of the statistical specification rather than as secondary technical details.
